\section{Síntesis y ampliación}

\subsection{Puntos clave de esta clase}

\begin{enumerate}
    \item \textbf{El problema del vacío de datos}: existen conjuntos de datos de imágenes a escala de miles de millones, pero la cantidad de datos para otras modalidades como modelos 3D es mucho menor. Cómo utilizar el conocimiento previo en 2D para compensar la falta de datos en 3D constituye la motivación central.

    \item \textbf{Neural Rendering como puente}: la función de renderizado diferenciable $g(\theta, \pi)$ mapea una representación 3D a una imagen 2D, permitiendo que las señales de optimización del espacio 2D se propaguen hacia atrás a los parámetros 3D.

    \item \textbf{De CLIP a los modelos difusos}: CLIP proporciona alineación a nivel semántico, pero carece de detalles. Los modelos difusos contienen un conocimiento previo mucho más rico a nivel de píxel; SDS ofrece una manera efectiva de aprovechar estos conocimientos previos.

    \item \textbf{El mecanismo central de SDS}: utiliza la función de pérdida de predicción de ruido del modelo difuso preentrenado como medida de calidad de renderizado, logrando una optimización eficiente mediante una simplificación que ignora el Jacobiano de U-Net.

    \item \textbf{Éxito y limitaciones de DreamFusion}: demuestra el enorme potencial de SDS en la generación texto a 3D, pero problemas como Janus Problem revelan las limitaciones inherentes del conocimiento previo en 2D.

    \item \textbf{Tendencias futuras}: la línea entre los modelos autoregresivos y difusos se está desdibujando; los marcos de generación multimodal unificados son una dirección importante.
\end{enumerate}

\subsection{Tabla comparativa de métodos}

\begin{table}[H]
\centering
\begin{tabular}{p{0.18\textwidth} p{0.22\textwidth} p{0.24\textwidth} p{0.22\textwidth}}
\toprule
Método & Ventaja central & Limitación principal & Escenario de uso \\
\midrule
CLIP guiado 3D & Alineación semántica directa e implementación relativamente simple & Detalles débiles, deriva de forma fácil & Exploración de conceptos sin ejemplos \\
SDS / DreamFusion & Permite reutilizar un conocimiento previo en difusos 2D fuerte & Problema Janus, sobre-saturación, inestabilidad geométrica & Generación de prototipos de alta fidelidad texto a 3D \\
VSD / distilación mejorada & Mejor diversidad y estabilidad & Entrenamiento y optimización más complejos & Generación 3D buscando mayor calidad y menos artefactos \\
Marco de generación multimodal unificado & Oportunidad de conectar las fronteras entre 2D/3D/video & Aún en etapa rápida de evolución & Investigación de modelos generativos universales de próxima generación \\
\bottomrule
\end{tabular}
\caption{Comparación de varias categorías de métodos centrales presentados en la Clase 13}
\end{table}

\subsection{Preguntas para reflexionar más}

\begin{itemize}
    \item ¿Hasta qué punto puede sustituir el conocimiento previo en 2D a los datos reales 3D sin quedar permanentemente limitado por problemas de consistencia de perspectiva?
    \item Cuando la línea entre los modelos autoregresivos y difusos se desdibuja continuamente, ¿los sistemas futuros de generación 3D conservarán las categorías de métodos claramente definidas de hoy?
    \item En el futuro, ¿Neural Rendering jugará más un papel de puente para el entrenamiento o se convertirá en una representación intermedia estándar dentro de la generación multimodal unificada?
\end{itemize}

\subsection{Lecturas adicionales}

\begin{itemize}
    \item Poole, B. et al. ``DreamFusion: Text-to-3D using 2D Diffusion.'' \textit{ICLR 2023}. \\
    ——Trabajo pionero que propone por primera vez el uso de SDS para la generación texto a 3D

    \item Wang, Z. et al. ``ProlificDreamer: High-Fidelity and Diverse Text-to-3D Generation with Variational Score Distillation.'' \textit{NeurIPS 2023}. \\
    ——Propone VSD (Distilación de Puntuaciones Variacional), mejorando significativamente los problemas de sobre-saturación y diversidad de SDS

    \item Mildenhall, B. et al. ``NeRF: Representing Scenes as Neural Radiance Fields for View Synthesis.'' \textit{ECCV 2020}. \\
    ——El artículo original de Neural Radiance Fields, un hito en el campo del Neural Rendering

    \item Kerbl, B. et al. ``3D Gaussian Splatting for Real-Time Radiance Field Rendering.'' \textit{SIGGRAPH 2023}. \\
    ——3D Gaussian Splatting, un método de renderizado diferenciable eficiente

    \item Jain, A. et al. ``DreamFields: Zero-Shot Text-Guided Object Generation with Dream Fields.'' \textit{CVPR 2022}. \\
    ——Generación 3D sin ejemplos guiada por CLIP

    \item Radford, A. et al. ``Learning Transferable Visual Models From Natural Language Supervision.'' \textit{ICML 2021}. \\
    ——Artículo original del modelo CLIP

    \item Luo, T. et al. ``Discrete Diffusion Modeling by Estimating the Ratios of the Data Distribution.'' \textit{ICML 2024}. \\
    ——Trabajo representativo de los modelos difusos discretos
\end{itemize}

%% ==================== Fin del cuerpo ====================

\end{document}
